Ifigenia
%Hablando de héroes malvados, Odiseo se lleva la palma y ese es el argumento central de mi \emph{Abandonando Ítaca}\footnote{\url{https://eolasediciones.es/libro/abandonando-itaca-leaving-ithaca/}}. De hecho, Ulises tiene las manos casi tan manchadas de sangre como Agamenón en el asesinato de la pobre Ifigenia.



Lo que nos lleva a Clitemnestra, la media hermana de Helena\footnote{Las dos son hijas de Leda, pero sus padres son distintos. Helena es hija de Zeus, Clitemnestra de Tindáreo. Hacerlas gemelas, como quiere Nolan, es una licencia audaz que permite que la divinal Lupita Nyong'o las interprete a ambas.}. La escena que nos presenta Nolan, en la que Clitemnestra asesina a Agamenón, no aparece en Homero, sino en la Orestíada (Eurípides). En la \emph{Odisea}, el relato de la muerte del rey, se ofrece en varios fragmentos. Casi al principio de la obra, Zeus explica que Egisto seduce a la esposa de Agamenón y lo mata a él, pese a la advertencia divina, recibiendo el castigo correspondiente por medio de Orestes. 


Fue el Padre de hombres y dioses primero en hablar entre ellos,
pues de su pecho le vino de Egisto preclaro el recuerdo
a quien el bien oído Orestes, el de Agamenón, dejó muerto.		
Y en recordándolo a él sus palabras les dio a los eternos:

Ay, los humanos siempre andan echando a los dioses la culpa;
que es que les van de nosotros sus males repiten, y nunca
ven que no estaba en su sino el revés por sus propias locuras.
Ahí está Egisto, que allende su sino quedaban las nupcias			
con  la mujer del Atrida, y matarlo al volver de la lucha. 
Fue —y lo sabía— su fin: de nosotros tenía segura
hecha advertencia por Hermes, vigía Argifontes, ‘escucha:  
que ni lo mates a él, ni andes tú de su esposa a la busca,
pues que de Orestes vendrá el talión del Atrida, cuando una		
fuerte añoranza le llegue, en flor de la edad, de la cuna’.
Hermes habló, mas con toda su buena intención la cordura
no trajo a Egisto, y ahora las paga por fin todas juntas.

El discurso de Zeus no tiene desperdicio. De entrada, rechaza la responsabilidad de los dioses en asuntos humanos: <<Ay, los humanos siempre andan echando a los dioses la culpa ;/ que es que les van de nosotros sus males repiten, y nunca /
ven que no estaba en su sino el revés por sus propias locuras.>> No anda desencaminado por una parte, pero por otra no faltan ejemplos a raudales, de todo lo contrario. De hecho, a lo largo de la Ilíada y la Odisea, los dioses no paran de enredar. Aquí, una vez más, la ambigüedad del Poeta es grande y misteriosa. Quizás todos esos enredos, parece decirnos, no son sino proyecciones de nuestras propias miserias. 

En todo caso, Zeus le echa la culpa exclusivamente a Egisto. A Clitemnestra, ni la nombra. Más tarde, durante la visita de Telémaco de Nestor, en Pylos, buscando noticias de su padre, el muchacho le pregunta al anciano rey por la suerte de Agamenón. 

Ah Néstor el de Neleo, verdad di tú verdadera:
¿cómo murió Agamenón, Atrida de muy larga alteza?
¿Y Menelao dónde andaba? ¿Qué traza de muerte le inventa
el urdetretas Egisto, que mata a quien más grande era?			


Y Nestor, en su respuesta, explica que el traidor consigue seducir a la esposa del caudillo ausente.

Que ella al principio decía que no a esos trajines obscenos,		
la divinal Clitemnestra, señora en sus sentimientos;
la acompañaba un cantor, que el Atrida a uno dio mandamiento
de bien guardarle a su esposa, a Troya apenas partiendo.
Mas cuando a ser ya domada la ató el destino del cielo,
él, arrastrando al cantor hasta un secano desierto,				
lo abandonó por de pájaros ser la presa y obsequio;
y se la sube a su casa con un quiero sí y un sí quiero.

Ya en la corte de Menelao, durante la cena, este cuenta como Egisto asesina a Agamenón, sin que Clitemnestra intervenga:




%\footnote{No olvidemos que Agamenón sacrifica a la hija de ambos, Ifigenia, degollándola él mismo para que los dioses les concedan el viento que las naves necesitan para partir a Troya.}. 

\begin{verse}
De una atalaya lo vio el atalayero que había\\
puesto allí Egisto urdetretas, con paga por él prometida\\
de dos talentos de oro; y un año iba ya de vigía,\\
no fuera a entrar sin ser visto quien fuerza y brío no olvida.\\
Fue a la carrera al pastor de los pueblos por darle noticia.\\
Pensando Egisto ya estaba su malartimaña zaína:\\
por la ciudad veinte hombres tomó, los de más osadía,\\
y colocó la emboscada al pie del banquete que haría.\\
A Agamenón, pastor de los pueblos, él fue en bienvenida\\
con los caballos y carros, rumiando maldades indignas.\\
Y lo mató, tras alzarlo, sin darse aquel cuenta, a la ruina,\\
mientras cenaba, como uno que a buey muerte da en boyeriza.\\
Y del Atrida ni un hombre quedó de los que le seguían,\\
ni uno de Egisto, que allí entre sus salas dejaron la vida.
\end{verse}

%No desperdicia epítetos para Egisto el poeta. Lo llama «urdetretas». Y por supuesto es un cobarde. Y por supuesto, Clitemnestra está compinchada con él, pero no es la ejecutora del crimen\footnote{Nolan nos mete aquí gato por liebre, o mejor dicho Esquilo por Homero. Las escenas de Nolan están tomadas de la \emph{Orestíada}, que se escribe en el 458 a.C., casi trescientos años más tarde que la \emph{Odisea}.}. Eso no impide que su hijo Orestes la mate en venganza.

La injusticia que Homero nos presenta cariperrunamente, o sea descaradamente, es flagrante. A la pobre Clitemnestra, el mostrenco de Agamenón (el personaje más odioso de la \emph{Ilíada}, de cuya oscuridad Nolan se hace eco vistiéndolo de Darth Vader) le mata a su hija del alma, una adolescente adorable, para conseguir el favor de los dioses (tiene narices, por cierto, que la Atenea de Nolan lo justifique, alegando que el valor del sacrificio es proporcional a lo que le cuesta al sacrificador, como si eso sirviera para explicar semejante salvajada). A continuación, el filicida se larga a asolar Troya durante diez años. Que Clitemnestra se echara un amante no solo parece justificable: se diría inevitable. Que se vengara del marido a su vuelta puede ser reprobable. ¿Pero hacía falta que su hijo la matara por eso? Por comparación, Menelao no solo no mata a Helena, ni le hace corte alguno en el rostro, sino que además, gracias a ella se agencia un pasaje en primera al Elíseo (escapándose del Hades). Se lo cuenta Proteo, el viejo del mar, en el mismo libro IV que estamos visitando.
